% The opening is computation first, followed by its geometric interpretation.
% Checked declarations are attached by the reader's rational-geometry pass.
Area, length, and volume can be computed before defining an integral of a
function. We begin with a concrete rational interval computation. We then
establish the elementary rules of area and identify the computation with the
quarter-circle. Throughout this chapter, the order of vertices determines
orientation; area and volume retain that orientation.

\section{A quarter-circle computation}
\label{thm:rational-circle-pythagorean}
Put
\[
 P(u)=\left(\frac{1-u^2}{1+u^2},\frac{2u}{1+u^2}\right),
 \qquad u_j=\frac{j}{2^n}\quad(0\le j\le2^n).
\]
All these coordinates are rational, and expansion gives
$P_1(u)^2+P_2(u)^2=1$. The endpoints are $(1,0)$ and $(0,1)$.
The chords and endpoint tangents in the picture motivate the two expressions
\[
 \ell(u,v)=\frac{(v-u)(1+uv)}{(1+u^2)(1+v^2)},
 \qquad h(u,v)=\frac{v-u}{1+uv}.
\]
For now, regard them simply as rational formulas. Define a raw computation by
\[
 Q_n=\left[
   \sum_{j=0}^{2^n-1}\ell(u_j,u_{j+1}),
   \sum_{j=0}^{2^n-1}h(u_j,u_{j+1})
 \right].
\]
Its first stages are $Q_0=[1/2,1]$ and $Q_1=[7/10,5/6]$.
The existing midpoint-update program computes exactly these sums; no new
schedule or numerical definition of $\pi$ is being substituted.

\paragraph{Validity of the computation.}\label{thm:quarter-circle-raw}
Finite refinement increases the lower sum and decreases the upper sum.
Furthermore
\[
 h(u,v)-\ell(u,v)=
 \frac{(v-u)^3}{(1+uv)(1+u^2)(1+v^2)}.
\]
On the chosen mesh each denominator is at least one. Summing the cell gaps
therefore gives
\[
 0\le Q_n^+-Q_n^-\le2^n(2^{-n})^3=4^{-n}.
\]
Thus the intervals are ordered, nested, and arbitrarily narrow: $Q$ is a
valid raw computable number. This conclusion uses finite rational formulas,
not an area or integral theorem.

\begin{figure}[htbp]\centering
\rationalCircleSubdivisionAnimation
\caption{A rational computation, whose geometric interpretation will follow.}
\label{fig:rational-parameter-line}
\end{figure}

\section{Area from ordered rational vectors}
\label{sec:rational-area}
A point or vector is a pair of rational numbers. The area of the
parallelogram spanned by the ordered vectors $u,v$ is
\[
 \omega(u,v)=u_1v_2-u_2v_1.
\]
This is the normalized alternating bilinear rule: it is additive and
rationally homogeneous in each input, vanishes when the two inputs agree,
and has $\omega((1,0),(0,1))=1$. Conversely those rules determine the
formula uniquely by expansion in the coordinate basis. We construct this
rule and prove its characterization; no new global axiom is introduced.

The area of the ordered triangle is
\[
 A[p,q,r]=\frac12\omega(q-p,r-p).
\]
A degenerate triangle has area zero. Exchanging two vertices reverses area;
a cyclic change of its three vertices preserves it. For every rational
point $o$, expansion gives the subdivision identity
\[
 A[p,q,r]=A[o,q,r]+A[p,o,r]+A[p,q,o].
\]
No condition that $o$ lies inside the triangle is needed: orientation handles
the cancellations.

For a rational affine map $x\mapsto Mx+b$, the transformation law is
\[
 A[Mp+b,Mq+b,Mr+b]=\det(M)A[p,q,r].
\]
Translations and proper rigid motions preserve area; reflections reverse it.
A homothety of ratio $s$ multiplies area by $s^2$. These assertions follow
from finite algebra, without lengths, angles, square roots, or trigonometry.

\section{Polygons and the choice of starting vertex}
\label{sec:rational-polygons}
A rational polygon is a cyclic sequence of rational points. In a stored list
$P=(p_0,\ldots,p_{N-1})$, the last vertex is joined back to the first;
there is no preferred first vertex and no assumption of simplicity or convexity.
Define
\[
 A(P)=\frac12\sum_{j=0}^{N-1}\omega(p_j,p_{j+1}),
 \qquad p_N=p_0.
\]
The empty list has area zero. Reversing the cycle reverses its area.

\paragraph{Theorem: a fan from any starting vertex.}
\label{thm:rational-polygon-fans}
For any rational point $o$,
\[
 \boxed{A(P)=\sum_{j=0}^{N-1}A[o,p_j,p_{j+1}].}
\]
In particular, take $o=p_0$. The first and last terms vanish, leaving
\[
 A(P)=\sum_{j=1}^{N-2}A[p_0,p_j,p_{j+1}].
\]
Choosing any other starting vertex gives the same result. More generally,
both the fan point and the place where the cyclic list starts may be changed
independently, without changing the computed rational number.

The proof is cancellation of the radial edges: each $\omega(o,p_j)$ occurs
once with each sign. The boundary terms remain. For a convex polygon this
is an ordinary triangulation. For other cycles it is an oriented triangle
sum; the triangles need not be disjoint or remain inside the polygon.
Self-crossing cycles therefore retain multiplicity rather than measuring
the union of the regions they trace. This distinction requires no change to
the formula or to the starting-vertex theorem.

\subsection{The same convention for volume}
For rational vectors in three dimensions put
$\Omega(u,v,w)=\det(u,v,w)$. The volume of the ordered tetrahedron is
\[
 V[p,q,r,s]=\frac16\Omega(q-p,r-p,s-p).
\]
Exchanging two vertices reverses volume, subdivision through any point is
additive with the displayed orientations, and a homothety of ratio $c$
multiplies volume by $c^3$. These computations precede the solid comparisons
later in this chapter. Spatial face areas and segment lengths may require
square roots; planar polygon area and tetrahedral volume do not.

\section{The geometric meaning of the computation}
\label{thm:geometric-pi-validity}
Return to the opening picture. The inner polygon has the origin followed by
the consecutive points $P(u_0),\ldots,P(u_{2^n})$. The outer polygon inserts
between consecutive circle points the intersection of their tangent lines
$P(u_j)\cdot z=1$ and $P(u_{j+1})\cdot z=1$.
These intersections are rational. Both boundary cycles have the positive
orientation shown in the picture.

\paragraph{Theorem: the interval endpoints are polygon areas.}
\label{thm:quarter-circle-polygons}
At every stage, not merely in the limit,
\[
 \boxed{Q_n=[A(P_n^-),A(P_n^+)].}
\]
Either polygon area can be computed from a fan at any starting vertex, or
indeed from any rational fan point. Its rational value is unchanged.
The proof identifies the triangle and tangent-cell formulas with the
already defined boundary-area computation.

The chords lie inside the quarter-disk and the tangent half-planes contain
it. Thus these are inscribed and circumscribed polygons. Their proved area
gap tends to zero. This is the area computation of the quarter-circle:
we have first constructed a number and then justified what its bounds measure.
No general area function on arbitrary regions is required.

Define $\pi=4Q$. The same construction on $[0,t]$ gives sector area $A(t)$,
so $Q=A(1)$. Finite polygon bounds give $2<\pi<4$, and rational coordinate
scaling gives the disk-area computation $\pi R^2$. For later use,
\[
 \frac{v-u}{1+v^2}\le A(v)-A(u)\le\frac{v-u}{1+u^2}
 \qquad(0\le u<v\le1).
\]
In particular $(v-u)/2\le A(v)-A(u)\le v-u$. These are geometric bounds;
the later integral constructions are unchanged and are not being redefined
as areas in this revision.

The sphere and cylinder arguments below are coordinate reconstructions in
the spirit of Archimedes' exhaustion, not transcriptions of his proofs.
The historical references remain at the end of the chapter.
